% !TEX root = ../ms-thesis.tex

\begin{thebibliography}{99}
  \bibitem{ref1} 吴丽, 梁皓, 虞华君, 等. 中国文化和旅游融合发展空间分异及驱动因素[J]. 经济地理, 2021, 41(2): 214-221. DOI:10.15957/j.cnki.jjdl.2021.02.023.

  \bibitem{ref2} 章坤, 朱海, 谢朝武. 中国文化和旅游产业高质量融合发展的适配关系与政策启示[J]. 自然资源学报, 2025, 40(4): 1084-1106. DOI:10.31497/zrzyxb.20250413.

  \bibitem{ref3} 郭艳萍, 刘敏. 基于 POI 数据的山西省旅游景区分类及空间分布特征[J]. 地理科学, 2021, 41(7): 1246-1255. DOI:10.13249/j.cnki.sgs.2021.07.015.

  \bibitem{ref4} 田志馥, 于亚娟, 黄辰宇. 基于 POI 数据挖掘的内蒙古县域旅游要素空间分布特征及影响因素[J]. 干旱区地理, 2025, 48(7): 1255-1266. DOI:10.12118/j.issn.1000-6060.2024.601.

  \bibitem{ref5} 王邵军, 李晓冰. 黄河流域高质量发展目标下非物质文化遗产与旅游业耦合协调发展研究：以山东省沿黄九市为例[J]. 东岳论丛, 2023, 44(11): 132-147. DOI:10.15981/j.cnki.dongyueluncong.2023.11.015.

  \bibitem{ref6} 王公为. 文旅融合背景下非物质文化遗产与旅游业的耦合发展研究：以内蒙古为例[J/OL]. (2020-03)[2026-04-14]. \url{http://ycfytsk.jx.chaoxing.com/info/2611}.

  \bibitem{ref7} FENSEL A, AKBAR Z, KARLE E, et al. Knowledge Graphs for Online Marketing and Sales of Touristic Services[J]. Information, 2020, 11(5): 253. DOI:10.3390/info11050253.

  \bibitem{ref8} SONG S, YANG C, XU L, et al. TravelRAG: A Tourist Attraction Retrieval Framework Based on Multi-Layer Knowledge Graph[J]. ISPRS International Journal of Geo-Information, 2024, 13(11): 414. DOI:10.3390/ijgi13110414.

  \bibitem{ref9} CADEDDU A, CHESSA A, DE LEO V, et al. Optimizing Tourism Accommodation Offers by Integrating Language Models and Knowledge Graph Technologies[J]. Information, 2024, 15(7): 398. DOI:10.3390/info15070398.

  \bibitem{ref10} XIAO D, WANG N, YU J, et al. A Practice of Tourism Knowledge Graph Construction based on Heterogeneous Information[C]//Proceedings of the 19th China National Conference on Computational Linguistics (CCL 2020). Haikou: Chinese Information Processing Society of China, 2020: 939-949. \url{https://aclanthology.org/2020.ccl-1.87.pdf}.

  \bibitem{ref11} 翁钢民, 李凌雁. 中国旅游与文化产业融合发展的耦合协调度及空间相关分析[J/OL]. 经济地理, 2016[2026-05-09]. \url{https://www.jjdl.com.cn/CN/abstract/article/1000-8462/32296}.

  \bibitem{ref12} 王兆峰, 谢佳亮. 中国文化和旅游融合发展效率时空动态演化及其驱动机制[J/OL]. 旅游学刊, 2024[2026-05-09]. \url{http://www.qikan.com.cn/article/lyxk20240111.html}.

  \bibitem{ref13} 何静, 冯学钢. 中国省际文旅融合空间关联网络结构演化及形成机制[J/OL]. 地理研究, 2025[2026-05-09]. \url{https://www.dlyj.ac.cn/CN/10.11821/dlyj020241213}.

  \bibitem{ref14} 朱媛媛, 等. 长江中游城市群“文—旅”产业融合发展的空间效应及驱动机制研究[J/OL]. 地理科学进展, 2022[2026-05-09]. \url{https://www.progressingeography.com/CN/10.18306/dlkxjz.2022.05.004}.

  \bibitem{ref15} 周海霞, 等. 长征沿线文化—旅游系统适配驱动文旅融合发展的时空演变与机理[J/OL]. 地理科学进展, 2025[2026-05-09]. \url{https://www.progressingeography.com/CN/10.18306/dlkxjz.2025.09.009}.

  \bibitem{ref16} 陆文镔, 等. 数智时代文旅融合的耦合机理研究：基于数字产业与旅游产业协同集聚视角[J/OL]. 自然资源学报, 2025[2026-05-09]. \url{https://www.jnr.ac.cn/CN/10.31497/zrzyxb.20251207}.

  \bibitem{ref17} 许春晓, 等. 文旅深度融合与资源创新利用：热现象与冷思考[J/OL]. 自然资源学报, 2025[2026-05-09]. \url{https://www.jnr.ac.cn/CN/top_access}.

  \bibitem{ref18} 张宏磊, 等. 中国现代旅游业体系建设中的旅游资源创新开发：理论认知与应用创新[J/OL]. 自然资源学报, 2025[2026-05-09]. \url{https://www.jnr.ac.cn/CN/top_access}.
\end{thebibliography}
